\documentclass[10pt,a4paper]{amsart}
\usepackage{amsmath,amssymb,amsthm,graphicx}
% For commutative diagrams we use Xy-pic
\usepackage[all]{xy}

% We use 2cell for 2-commutative diagrams.
\xyoption{2cell}
% Theorem environments.
%
\theoremstyle{plain}
\newtheorem{theorem}[subsection]{Theorem}
\newtheorem{proposition}[subsection]{Proposition}
\newtheorem{lemma}[subsection]{Lemma}

\theoremstyle{definition}
\newtheorem{definition}[subsection]{Definition}
\newtheorem{example}[subsection]{Example}
\newtheorem{exercise}[subsection]{Exercise}
\newtheorem{situation}[subsection]{Situation}

\theoremstyle{remark}
\newtheorem{remark}[subsection]{Remark}
\newtheorem{remarks}[subsection]{Remarks}

\numberwithin{equation}{subsection}

% Macros
%
\def\lim{\mathop{\mathrm{lim}}\nolimits}
\def\colim{\mathop{\mathrm{colim}}\nolimits}
\def\Spec{\mathop{\mathrm{Spec}}}
\def\Hom{\mathop{\mathrm{Hom}}\nolimits}
\def\Ext{\mathop{\mathrm{Ext}}\nolimits}
\def\SheafHom{\mathop{\mathcal{H}\!\mathit{om}}\nolimits}
\def\SheafExt{\mathop{\mathcal{E}\!\mathit{xt}}\nolimits}
\def\Sch{\mathit{Sch}}
\def\Mor{\mathop{\mathrm{Mor}}\nolimits}
\def\Ob{\mathop{\mathrm{Ob}}\nolimits}
\def\Sh{\mathop{\mathit{Sh}}\nolimits}
\def\NL{\mathop{N\!L}\nolimits}
\def\CH{\mathop{\mathrm{CH}}\nolimits}
\def\proetale{{pro\text{-}\acute{e}tale}}
\def\etale{{\acute{e}tale}}
\def\QCoh{\mathit{QCoh}}
\def\Ker{\mathop{\mathrm{Ker}}}
\def\Im{\mathop{\mathrm{Im}}}
\def\Coker{\mathop{\mathrm{Coker}}}
\def\Coim{\mathop{\mathrm{Coim}}}

% Boxtimes
%
\DeclareMathSymbol{\boxtimes}{\mathbin}{AMSa}{"02}

%
% Macros for moduli stacks/spaces
%
\def\QCohstack{\mathcal{QC}\!\mathit{oh}}
\def\Cohstack{\mathcal{C}\!\mathit{oh}}
\def\Spacesstack{\mathcal{S}\!\mathit{paces}}
\def\Quotfunctor{\mathrm{Quot}}
\def\Hilbfunctor{\mathrm{Hilb}}
\def\Curvesstack{\mathcal{C}\!\mathit{urves}}
\def\Polarizedstack{\mathcal{P}\!\mathit{olarized}}
\def\Complexesstack{\mathcal{C}\!\mathit{omplexes}}
% \Pic is the operator that assigns to X its picard group, usage \Pic(X)
% \Picardstack_{X/B} denotes the Picard stack of X over B
% \Picardfunctor_{X/B} denotes the Picard functor of X over B
\def\Pic{\mathop{\mathrm{Pic}}\nolimits}
\def\Picardstack{\mathcal{P}\!\mathit{ic}}
\def\Picardfunctor{\mathrm{Pic}}
\def\Deformationcategory{\mathcal{D}\!\mathit{ef}}
\usepackage{cleveref}
\crefname{originalreference}{original reference}{original references}
\Crefname{originalreference}{Original reference}{Original references}
\makeatletter
\newlabel{lemma-glue-etale-sheaf-check-after-base-change}{{0GF1}{source}{}{sourcepointer-1}{}}
\newlabel{lemma-glue-etale-sheaf-check-after-base-change@cref}{{[originalreference][0][]0GF1}{[1][0][]source}{}{}{}}
\makeatother
\crefname{originalreference}{original reference}{original references}
\Crefname{originalreference}{Original reference}{Original references}
\makeatletter
\newlabel{lemma-glue-etale-sheaf-integral-surjective}{{0GEZ}{source}{}{sourcepointer-2}{}}
\newlabel{lemma-glue-etale-sheaf-integral-surjective@cref}{{[originalreference][0][]0GEZ}{[1][0][]source}{}{}{}}
\makeatother
\crefname{originalreference}{original reference}{original references}
\Crefname{originalreference}{Original reference}{Original references}
\makeatletter
\newlabel{lemma-glue-etale-sheaf-section}{{0GEY}{source}{}{sourcepointer-3}{}}
\newlabel{lemma-glue-etale-sheaf-section@cref}{{[originalreference][0][]0GEY}{[1][0][]source}{}{}{}}
\makeatother
\crefname{originallemma}{original lemma}{original lemmas}
\Crefname{originallemma}{Original lemma}{Original lemmas}
\makeatletter
\newlabel{more-morphisms-lemma-dominate-fppf-finite}{{0CP0}{source}{}{sourcepointer-4}{}}
\newlabel{more-morphisms-lemma-dominate-fppf-finite@cref}{{[originallemma][0][]0CP0}{[1][0][]source}{}{}{}}
\makeatother
\title{Etale sheaves satisfy descent for a surjective flat locally finitely presented morphism}
\author{The Stacks Project Authors}
\date{stacks-project:etale-cohomology.tex:local-snapshot-20261011; GNU FDL 1.2 or later}
\begin{document}
\maketitle
\noindent\textit{Editorial note.} All delivered mathematical text is the author's own, copied from the pinned Stacks Project chapter file etale-cohomology.tex; no statement or proof step was rewritten, reordered or removed.\par\medskip
\noindent\textit{Editorial note.} The counted proof is the author's complete argument: the proof is written out in full: the equalizer quasi-inverse, the reduction to f\textasciicircum{}\{-1\}G = F checked locally on Y, the finite cover dominating the fppf cover, the after-base-change reduction, and the final appeal to the sections case.\par\medskip
\noindent\textit{Editorial note.} The label lemma-glue-etale-sheaf-fppf-cover names the counted statement, and the proof environment immediately following it in the source is delivered verbatim as its proof.\par\medskip
\noindent\textit{Editorial note.} The counted statement is carried by the author's own lemma environment in the chapter source.\par\medskip
\noindent\textit{Editorial note.} The statement's descent datum notation is the author's and is delivered verbatim; the four cited inputs are kept as external pointers with their canonical Stacks tags.\par\medskip
\noindent\textit{Editorial note.} The preamble is the chapter's own native preamble excerpt already used for this source: the Xy-pic setup (native preamble lines 14-18) and the native theorem declarations with the mathematical macros (native preamble lines 154-217).\par\medskip
\noindent\textit{Allowed standard prerequisites (editorial).}\begin{itemize}
\item The equalizer description of the quasi-inverse for the integral case. Same chapter, the lemma establishing descent for a surjective integral morphism with the equalizer construction of the glued sheaf. Supplies the candidate quasi-inverse G = Equalizer(f\_*F => f\_\{1,*\}F\_1) and the reduction to proving f\textasciicircum{}\{-1\}G = F.
\item Checking the comparison after a base change. Same chapter, the lemma allowing the verification of the gluing statement to be performed after a base change by a finite morphism. Supplies the passage from Y to a finite cover Z -> Y and to Z x\_Y Z.
\item Domination of an fppf cover by a finite surjective morphism. More on Morphisms, the lemma producing a finite surjective morphism which Zariski locally factors through a given surjective flat locally finitely presented morphism. Supplies the cover Z -> Y and the open covering Z = union W\_i with W\_i -> Y factoring through X.
\item Descent in the case where the morphism has a section. Same chapter, the lemma treating the gluing problem for a morphism with a section. Supplies the final step once the morphism has sections Zariski locally.
\end{itemize}\medskip
\ifdefined\hypertarget\hypertarget{sourcepointer-1}{}\fi
\noindent\textit{Source reference.} Original reference 0GF1: sources/stacks-etale-cohomology/pinned-original-source.tex, source lines 23068--23068. Prerequisite pointer invoked at proof line(s) 23136; it names a standard fact of the same chapter that the counted argument uses, and every step of the counted proof built on it is retained verbatim in the delivered body.\par\medskip
\ifdefined\hypertarget\hypertarget{sourcepointer-2}{}\fi
\noindent\textit{Source reference.} Original reference 0GEZ: sources/stacks-etale-cohomology/pinned-original-source.tex, source lines 22972--22972. Prerequisite pointer invoked at proof line(s) 23114; it names a standard fact of the same chapter that the counted argument uses, and every step of the counted proof built on it is retained verbatim in the delivered body.\par\medskip
\ifdefined\hypertarget\hypertarget{sourcepointer-3}{}\fi
\noindent\textit{Source reference.} Original reference 0GEY: sources/stacks-etale-cohomology/pinned-original-source.tex, source lines 22952--22952. Prerequisite pointer invoked at proof line(s) 23142; it names a standard fact of the same chapter that the counted argument uses, and every step of the counted proof built on it is retained verbatim in the delivered body.\par\medskip
\ifdefined\hypertarget\hypertarget{sourcepointer-4}{}\fi
\noindent\textit{Source reference.} Original lemma 0CP0: sources/stacks-chow/tags.txt, source lines 15910--15910. Prerequisite pointer invoked at proof line(s) 23134; it names a standard fact of a sibling chapter that the counted argument uses, and every step of the counted proof built on it is retained verbatim in the delivered body.\par\medskip
\begin{lemma}
\label{lemma-glue-etale-sheaf-fppf-cover}
Let $f : X \to Y$ be a morphism of schemes which is
surjective, flat, locally of finite presentation.
The functor
$$
\Sh(Y_\etale)
\longrightarrow
\text{descent data for \'etale sheaves wrt }\{X \to Y\}
$$
is an equivalence of categories.
\end{lemma}

\begin{proof}
Exactly as in the proof of
Lemma \ref{lemma-glue-etale-sheaf-integral-surjective}
we claim a quasi-inverse is given by the functor sending
$(\mathcal{F}, \varphi)$ to
$$
\mathcal{G} =
\text{Equalizer}\left(
\xymatrix{
f_*\mathcal{F}
\ar@<1ex>[r] \ar@<-1ex>[r] &
f_{1, *}\mathcal{F}_1
}
\right)
$$
and in order to prove this it suffices to show that
$f^{-1}\mathcal{G} \to \mathcal{F}$ is an isomorphism.
This we may check locally, hence we may and do assume $Y$
is affine. Then we can find a finite surjective morphism
$Z \to Y$ such that there exists an open covering
$Z = \bigcup W_i$ such that $W_i \to Y$ factors through $X$.
See
More on Morphisms, Lemma \ref{more-morphisms-lemma-dominate-fppf-finite}.
Applying Lemma
\ref{lemma-glue-etale-sheaf-check-after-base-change}
we see that it suffices to prove
the lemma after replacing $Y$ by $Z$ and $Z \times_Y Z$ and $f$
by its base change. Thus we may assume $f$ has sections Zariski locally.
Of course, using that the problem is local on $Y$ we reduce
to the case where we have a section which is
Lemma \ref{lemma-glue-etale-sheaf-section}.
\end{proof}

\end{document}
